\section{Introduction}

Les tumeurs cérébrales comptent parmi les pathologies neurologiques les plus graves.
Leur classification précise et rapide en gliome, méningiome et tumeur hypophysaire
est indispensable pour orienter les décisions chirurgicales et thérapeutiques~\cite{ref_bioengineering11030266}.
L'imagerie par résonance magnétique (IRM) constitue la modalité diagnostique de référence,
mais l'interprétation manuelle par les radiologues reste chronophage et sujette à une
variabilité inter-observateur.

L'apprentissage automatique offre une voie prometteuse vers des pipelines de classification
automatisés et reproductibles. Si les modèles d'apprentissage profond (CNN) rapportent
des performances élevées sur de grands jeux de données~\cite{ref_diagnostics15212692},
le ML classique avec descripteurs artisanaux reste compétitif sur des corpus modestes
et, surtout, offre une interprétabilité : les caractéristiques utilisées pour la
classification sont explicitement compréhensibles par l'opérateur.

\subsection{Périmètre et motivation}

Ce travail se concentre exclusivement sur le ML classique avec descripteurs artisanaux.
Nous ne détaillons pas les dérivations mathématiques des algorithmes ; l'accent est
mis sur la méthodologie expérimentale et les résultats obtenus. Les questions principales
sont les suivantes :

\begin{itemize}
  \item Quels paramètres de descripteurs séparent le mieux les trois classes tumorales ?
  \item La fusion de plusieurs familles de caractéristiques améliore-t-elle une famille isolée ?
  \item Quel classifieur tire le plus de profit des caractéristiques optimisées ?
  \item Les résultats sont-ils reproductibles et statistiquement significatifs ?
\end{itemize}

\subsection{Contributions}

\begin{enumerate}
  \item Un cadre d'optimisation de paramètres guidé par la séparabilité pour des
        descripteurs IRM artisanaux, explorant \NPhaseOneConfigs{} configurations
        sans entraînement de classifieur.
  \item Un benchmark complet de \NModels{} classifieurs sur \NFeatureSets{} jeux
        de caractéristiques, soit \NRuns{} pipelines évalués.
  \item Une validation statistique multicouche (validation croisée, intervalles de
        confiance bootstrap, test de McNemar, test de Friedman, stabilité par graine).
  \item L'identification du meilleur pipeline : \BestPipeline{} avec \BestAcc{} de précision.
\end{enumerate}
